\chapter{Characters, Euler Products and the Sieve}\label{ch:characters}

\begin{center}\itshape
Euler factored the harmonic series into a product over primes.\\
The same idea, with a sign that remembers a residue class, is the door to Dirichlet's theorem.
\end{center}


An integer may have many properties, but a sum often needs a simple numerical signal. Define a rule on integers that returns $0$ for even numbers, $1$ for numbers congruent to $1$ modulo $4$, and $-1$ for numbers congruent to $3$ modulo $4$. This is the character $\chi_4$. It remembers a residue class and is multiplicative. The first four values, at $1,2,3,4$, are $1,0,-1,0$; the next four repeat them.

Why is this useful? A sum weighted by $\chi_4$ distinguishes the two odd residue classes without listing their members separately. A product of such functions keeps their arithmetic structure. When we combine a character with a series of the form $\sum_{n\ge1}\chi(n)n^{-s}$, multiplication over integers becomes an Euler product over primes, provided the series converges absolutely.

This chapter joins two earlier subjects. Primitive roots and finite unit groups explain which multiplicative signs or roots of unity a character can take. The Möbius function explains how to remove integers divisible by small primes. The resulting sieve is an exact finite count before any limit or asymptotic estimate enters.

An Euler product is a bridge to the prime number theorem and Dirichlet's theorem on primes in residue classes. It is not itself a proof of either theorem. We will be explicit about the convergence conditions and about the point where deeper analysis is required.


\section{Dirichlet characters}

% t1-pass2 motive 81f967b9
Remember the Legendre symbol: it is zero on multiples of $p$, and on the other residue classes it tells us which are squares. A Dirichlet character keeps the same pattern in a wider setting. It assigns a root of unity to each invertible residue class, respects multiplication, and vanishes on the other classes. The benefit is not just a new definition. Summing characters lets us isolate one residue class by cancellation, exactly as roots of unity filtered coefficients in Volume~1.

\begin{definition}\label{def:t1-character}
Let $m\ge1$. A \emph{Dirichlet character modulo $m$} is a function $\chi:\Z\to\C$ such that
\begin{enumerate}[label=(\roman*)]
\item $\chi(n)=0$ whenever $\gcd(n,m)>1$, and $\chi(n)\ne0$ otherwise;
\item $\chi(n+m)=\chi(n)$ for every $n$;
\item $\chi(ab)=\chi(a)\chi(b)$ for all integers $a,b$.
\end{enumerate}
The \emph{principal character} $\chi_0$ is defined by $\chi_0(n)=1$ if $\gcd(n,m)=1$ and $\chi_0(n)=0$ otherwise.
\end{definition}

Condition~(iii) on the units is a homomorphism $(\Z/m\Z)^\times\to\C^\times$, and condition~(i) is the only way to extend it by zero while keeping complete multiplicativity. Conversely, every homomorphism of the unit group extends in this way: if a prime dividing $m$ divides $ab$, it divides $a$ or $b$, so both sides of~(iii) vanish off the units. Thus Dirichlet characters modulo $m$ are precisely the homomorphisms $(\Z/m\Z)^\times\to\C^\times$, extended by zero. There is nothing further to choose.

\begin{examplex}[three characters one already knows]
The principal character modulo $m$ is one. The function $\chi_4$ of Theorem~\ref{thm:r2},
\[
\chi_4(n)=\begin{cases}
0&n\text{ even},\\
1&n\equiv1\pmod4,\\
-1&n\equiv3\pmod4,
\end{cases}
\]
is the non-principal character modulo $4$. The Legendre symbol $n\mapsto\bigl(\frac np\bigr)$, for an odd prime $p$, is a character modulo $p$: it is periodic, vanishes at multiples of $p$, and is completely multiplicative. For the prime $2$, the supplementary law is itself a character of modulus $8$. Define
\[
\chi_8(n)=\begin{cases}
0&n\text{ even},\\
1&n\equiv\pm1\pmod8,\\
-1&n\equiv\pm3\pmod8.
\end{cases}
\]
On the units $\{1,3,5,7\}$ one has $\chi_8(1)=1$, $\chi_8(3)=\chi_8(5)=-1$ and $\chi_8(7)=1$. Each non-identity unit has order $2$, and the products $3\cdot5\equiv7$, $3\cdot7\equiv5$ and $5\cdot7\equiv3$ match the products of the values. So $\chi_8$ is a character modulo $8$. Corollary~\ref{cor:two} says $\bigl(\frac2p\bigr)=\chi_8(p)$ for every odd prime $p$.
\end{examplex}

\begin{figure}[ht]
\centering
\fig{ch14-chi8}
\caption{The eight values of $\chi_8$. Even residues contribute $0$. The two classes $\pm1$ receive $+1$, and the two classes $\pm3$ receive $-1$. Reading the clock is the supplementary law for the symbol $\bigl(\frac2{\cdot}\bigr)$.}
\label{fig:t1-chi8}
\end{figure}

\begin{proposition}\label{prop:t1-charsum}
Let $\chi$ be a character modulo $m$. Then
\[
\sum_{n=1}^m\chi(n)=\begin{cases}\varphi(m)&\text{if }\chi=\chi_0,\\0&\text{otherwise.}\end{cases}
\]
\end{proposition}
\begin{proof}
A character vanishes on non-units, so the sum
over all residue classes is the sum over
$(\Z/m\Z)^\times$.
For the principal character $\chi_0$, each
of the $\varphi(m)$ units contributes $1$.
The sum is therefore $\varphi(m)$.

If $\chi$ is not principal, choose a unit $b$
with $\chi(b)\ne1$. Multiplication by $b$
permutes the unit classes. Writing $S$ for
the sum of their character values, complete
multiplicativity gives
\[
S=\sum_{u\in(\Z/m\Z)^\times}\chi(bu)
 =\chi(b)\sum_{u\in(\Z/m\Z)^\times}\chi(u)
 =\chi(b)S.
\]
Hence $(\chi(b)-1)S=0$. The factor
$\chi(b)-1$ is nonzero by the choice of
$b$, so $S=0$. This is exact cancellation
over a finite group; no convergence theorem
is needed.
\end{proof}

The count of all such $\chi$ needs the shape of the unit group, which is now complete.

\begin{proposition}\label{prop:t1-unitproduct}
Let $m=2^k\prod_ip_i^{a_i}$ with odd primes $p_i$. The Chinese remainder theorem identifies
\[
(\Z/m\Z)^\times
\]
with the product of $(\Z/2^k\Z)^\times$ and the groups $(\Z/p_i^{a_i}\Z)^\times$. Each odd factor is cyclic, by Theorem~\ref{thm:cyclicn}. The factor of $2$ is cyclic for $k\le2$, and for $k\ge3$ it is the product described by Theorem~\ref{thm:twopower}. In particular $(\Z/m\Z)^\times$ is a product of cyclic groups, and it has $\varphi(m)$ elements.
\end{proposition}
\begin{proof}
Iterating Theorem~\ref{thm:crt} gives a ring isomorphism $\Z/m\Z\cong\prod\Z/p^{a}\Z$, the product running over the prime powers in $m$, and a ring isomorphism identifies units with tuples of units. The cyclicity statements are Theorem~\ref{thm:cyclicn} for the odd prime powers and for $2$ and $4$, and Theorem~\ref{thm:twopower} for higher powers of $2$.
\end{proof}

\begin{theorem}\label{thm:t1-charcount}
There are exactly $\varphi(m)$ Dirichlet characters modulo $m$. If $a$ is a unit and $a\not\equiv1\pmod m$, some character satisfies $\chi(a)\ne1$.
\end{theorem}
\begin{proof}
A cyclic group of order $N$, generated by $g$, has exactly $N$ homomorphisms to $\C^\times$: the image of $g$ may be any $N$th root of unity $e^{2\pi i j/N}$, $0\le j<N$, and each choice extends uniquely by $\chi(g^k)=e^{2\pi i jk/N}$. A homomorphism from a direct product of groups to $\C^\times$ is a tuple of homomorphisms, one from each factor, because $\chi(g,h)=\chi(g,1)\chi(1,h)$. Proposition~\ref{prop:t1-unitproduct} therefore supplies exactly $\varphi(m)$ characters.

For the second claim, write the unit group as a product of cyclic groups and let $a$ have a non-trivial component $g^k$ in a cyclic factor of order $N$, with $N\nmid k$. The character of that factor sending $g$ to $e^{2\pi i/N}$ does not fix $g^k$, and the trivial character on every other factor extends it to the whole unit group.
\end{proof}

\begin{think}
Modulo $5$, a primitive root is $2$, and the units are $2^0,2^1,2^2,2^3\equiv1,2,4,3$. How many characters should there be, and which of them is the Legendre symbol? Compute $\bigl(\frac2 5\bigr)$ from the supplementary law before looking at the example.
\end{think}

\begin{examplex}[the character table modulo $5$]
Since $\varphi(5)=4$, there are four characters. With generator $2$ they are $\chi_j(2^k)=e^{2\pi i jk/4}=i^{jk}$ for $j=0,1,2,3$. On the units $1,2,3,4$ one finds
\[
\begin{array}{c|cccc}
n&1&2&3&4\\\hline
\chi_0&1&1&1&1\\
\chi_1&1&i&-i&-1\\
\chi_2&1&-1&-1&1\\
\chi_3&1&-i&i&-1.
\end{array}
\]
The symbol $\bigl(\frac n5\bigr)$ equals $1, -1, -1, 1$ on these four units, because $\bigl(\frac25\bigr)=-1$ by the rule for $p\equiv\pm3\pmod8$ and $\bigl(\frac35\bigr)=\bigl(\frac53\bigr)=\bigl(\frac23\bigr)=-1$. It is exactly $\chi_2$, the unique non-principal character with values in $\{\pm1\}$. A cyclic group of even order has exactly one element of order $2$ in its dual, so there is no other candidate.
\end{examplex}

\begin{theorem}[orthogonality]\label{thm:t1-ortho}
Let $\chi,\psi$ be characters modulo $m$, and let $a$ be an integer.
\begin{enumerate}[label=(\roman*)]
\item $\displaystyle\sum_{n=1}^m\chi(n)\overline{\psi(n)}=\begin{cases}\varphi(m)&\text{if }\chi=\psi,\\0&\text{otherwise.}\end{cases}$
\item If $\gcd(a,m)=1$, then $\displaystyle\sum_{\chi}\chi(a)=\begin{cases}\varphi(m)&\text{if }a\equiv1\pmod m,\\0&\text{otherwise,}\end{cases}$ the sum running over all $\varphi(m)$ characters.
\item Consequently, for any integer $n$,
\[
\frac1{\varphi(m)}\sum_{\chi}\chi(n)\overline{\chi(a)}
\]
equals $1$ if $n\equiv a\pmod m$ and $\gcd(a,m)=1$, and equals $0$ otherwise.
\end{enumerate}
\end{theorem}
\begin{proof}
Values of $\chi$ on units are roots of unity, so $\overline{\chi(n)}=\chi(n)^{-1}$ whenever $\chi(n)\ne0$, and the function $\chi\overline{\psi}$ is the character $\chi\psi^{-1}$. It is principal if and only if $\chi=\psi$. Part~(i) is therefore Proposition~\ref{prop:t1-charsum}.

For~(ii), the characters form a group under multiplication: the product of two homomorphisms is a homomorphism, and Theorem~\ref{thm:t1-charcount} says there are $\varphi(m)$ of them, so every homomorphism appears. If $a\equiv1$, every term is $1$. If not, some $\psi$ has $\psi(a)\ne1$, and multiplication by $\psi$ permutes the characters, so the sum $S$ satisfies $\psi(a)S=S$.

For~(iii), if $\gcd(a,m)>1$ then $\overline{\chi(a)}=0$ for every $\chi$. If $\gcd(a,m)=1$, then $\overline{\chi(a)}=\chi(a)^{-1}=\chi(a^{-1})$, so the sum equals $\sum_{\chi}\chi(na^{-1})$. This vanishes when $\gcd(n,m)>1$. When $n$ is also a unit, part~(ii) returns $\varphi(m)$ if $na^{-1}\equiv1\pmod m$ and returns $0$ otherwise. Dividing by $\varphi(m)$ is the claim.
\end{proof}

Part~(iii) is the device that isolates a single residue class. Volume~6 writes the indicator of the primes $p\equiv a\pmod m$ by inserting this average into a sum over primes. The identity itself is finite and exact; no analysis has entered yet.

% t1-pass2 examples 81f967b9
\begin{examplex}[the non-principal character modulo $4$]
On odd $n$, $\chi_4(n)=(-1)^{(n-1)/2}$. Thus $\chi_4(1)=1$, $\chi_4(3)=-1$, $\chi_4(5)=1$, $\chi_4(7)=-1$, and $\chi_4$ vanishes on evens. The sum over a complete set of residues modulo $4$ is $0$, as Proposition~\ref{prop:t1-charsum} requires for a non-principal character. This is the character in Jacobi's formula for $r_2(n)$.
\end{examplex}

\begin{examplex}[two rows modulo $5$]
With primitive root $2$, the character $\chi_2(2^k)=(-1)^k$ takes the values $1,-1,1,-1$ on $1,2,4,3$. On the ordered units $1,2,3,4$ the values are $1,-1,-1,1$, which is the Legendre symbol. The sum of this row against the principal row $1,1,1,1$ is $1-1-1+1=0$.
\end{examplex}

\section{Euler products}

% t1-pass2 motive 78236f4c
A product of geometric series produces all positive integers if we choose one power of each prime. Unique factorisation tells us that each integer is produced exactly once. For a completely multiplicative function, this observation turns a Dirichlet series into a product over primes. Before we rearrange infinitely many terms, however, we must check absolute convergence. The region $s>1$ lets us make that step legally; the boundary $s=1$ is a different problem, not a harmless substitution.

A completely multiplicative function turns a sum over the integers into a product over the primes, because unique factorisation says that every $n$ is built from primes in exactly one way. The cost is a little convergence, which Volume~1's integral comparison supplies as soon as the exponent exceeds $1$.

\begin{lemma}\label{lem:t1-pseries}
Let $s>1$. The series $\sum_{n\ge1}n^{-s}$ converges, and for every real $X\ge1$,
\[
\sum_{n>X}n^{-s}\le\int_{\lfloor X\rfloor}^{\infty}t^{-s}\,\mathrm dt=\frac{\lfloor X\rfloor^{1-s}}{s-1}.
\]
\end{lemma}
\begin{proof}
The function $t^{-s}$ decreases. For $n\ge2$, $n^{-s}\le\int_{n-1}^n t^{-s}\,\mathrm dt$, so the tail from $n=\lfloor X\rfloor+1$ onwards is at most the integral from $\lfloor X\rfloor$ to infinity.
\end{proof}

\begin{lemma}\label{lem:t1-dseries}
Let $s>1$, and let $\sum a_nn^{-s}$ and $\sum b_nn^{-s}$ converge absolutely. Then
\[
\Bigl(\sum_{n\ge1}a_nn^{-s}\Bigr)\Bigl(\sum_{m\ge1}b_mm^{-s}\Bigr)
=\sum_{N\ge1}(a*b)(N)\,N^{-s}.
\]
\end{lemma}
\begin{proof}
The double sum $\sum_{n,m}|a_nb_m|(nm)^{-s}$ is the product of the two absolute series, hence finite. Absolute convergence lets one group the pairs $(n,m)$ by the product $N=nm$. The inner sum at $N$ is $(a*b)(N)N^{-s}$.
\end{proof}

\begin{theorem}[Euler product]\label{thm:t1-eulerprod}
Let $s>1$ and let $f$ be completely multiplicative, with $f(1)=1$ and $|f(n)|\le1$ for every $n$. Then
\[
\sum_{n\ge1}f(n)\,n^{-s}=\prod_p\bigl(1-f(p)\,p^{-s}\bigr)^{-1},
\]
and the product converges to a non-zero complex number.
\end{theorem}
\begin{proof}
Expanding a finite product of geometric series gives
\[
P_X:=\prod_{p\le X}\bigl(1-f(p)p^{-s}\bigr)^{-1}
=\sum_{n\in S_X}f(n)\,n^{-s},
\]
where $S_X$ is the set of positive integers whose prime factors are all $\le X$. Every $n\le X$ lies in $S_X$. The full series converges absolutely by comparison with Lemma~\ref{lem:t1-pseries}. An integer outside $S_X$ has a prime factor $>X$, hence exceeds $X$, so
\[
\Bigl|\sum_{n\ge1}f(n)n^{-s}-P_X\Bigr|
\le\sum_{n>X}n^{-s}\longrightarrow0.
\]
Thus the partial products converge to the series. The limit is not zero. The inequality $1+t\le e^t$ gives $\prod_p(1+p^{-s})\le\exp(\sum_pp^{-s})<\infty$, since $\sum_pp^{-s}$ converges by comparison with Lemma~\ref{lem:t1-pseries}. Each factor satisfies $|1-f(p)p^{-s}|\le1+p^{-s}$, and a finite product of non-zero terms is non-zero, so
\[
|P_X|=\Bigl|\prod_{p\le X}\bigl(1-f(p)p^{-s}\bigr)\Bigr|^{-1}
\ge\Bigl(\prod_p(1+p^{-s})\Bigr)^{-1}>0.
\]
The lower bound does not depend on $X$. An infinite product of non-zero terms may still vanish: $\prod_p(1-1/p)=0$ is the reason $\sum 1/n$ diverges. For $s>1$ the comparison forbids that vanishing.
\end{proof}

% t1-pass2 why
\paragraph{Why this works.}
For $s>1$, absolute convergence allows us to expand each Euler factor as a geometric series and multiply the series together. Unique prime factorisation says that every positive integer appears in exactly one way, giving the Dirichlet series. The tail of the integer sum is controlled by the convergent series $\sum n^{-s}$, so the rearrangement is justified. Do not move the argument unchanged to $s=1$: the comparison series diverges there, and a new idea is needed.

Two special cases carry the names used later. For $f=1$,
\[
\zeta(s):=\sum_{n\ge1}n^{-s}=\prod_p(1-p^{-s})^{-1}\qquad(s>1).
\]
For a character $\chi$,
\[
L(s,\chi):=\sum_{n\ge1}\chi(n)\,n^{-s}=\prod_p\bigl(1-\chi(p)p^{-s}\bigr)^{-1}\qquad(s>1).
\]
Both are non-zero for real $s>1$. The product for $\zeta$ is an equality of positive real numbers. The product for $L$ is an equality in $\C$.

\begin{corollary}\label{cor:t1-mu-zeta}
For $s>1$, $\displaystyle\sum_{n\ge1}\mu(n)\,n^{-s}=\frac1{\zeta(s)}=\prod_p(1-p^{-s})$.
\end{corollary}
\begin{proof}
Lemma~\ref{lem:t1-dseries} and Theorem~\ref{thm:mobius} give $\zeta(s)\sum\mu(n)n^{-s}=\sum\varepsilon(n)n^{-s}=1$. The Euler product for $\zeta$ then inverts term by term.
\end{proof}

\begin{examplex}[a product one can compare with $\pi^2/6$]
At $s=2$, the series is $1+\frac14+\frac19+\cdots=\pi^2/6$, the Basel sum from Volume~1. The partial sum through $n=30$ equals $1.61215\ldots$. The product over the primes $p\le30$ equals $1.63307\ldots$, and $\pi^2/6=1.64493\ldots$. The product is ahead of the partial sum because it already includes every integer whose prime factors are at most $30$, including those larger than $30$, such as $32$, $49$ and $77$. Both approximations approach the same limit from below; they are not two competing guesses.
\end{examplex}

\begin{think}
Corollary~\ref{cor:t1-mu-zeta} says that $\sum\mu(n)/n^s$ tends to zero as $s\downarrow1$, because $\zeta(s)$ grows like the harmonic series. Does the single series $\sum\mu(n)/n$ converge? The Euler product at $s=1$ would be $\prod(1-1/p)=0$, which is a hint and not yet a proof: the identity was stated for $s>1$. Volume~6 returns to this series. What, in the argument above, used $s>1$ in an essential way?
\end{think}

The answer to the last question is absolute convergence. At $s=1$ the harmonic series diverges, the comparison that identified the sum with the product stops, and a conditionally convergent series cannot be rearranged with the same freedom. That is a precise description of the boundary, not a vague warning.

% t1-pass2 examples 78236f4c
\begin{examplex}[the product for $\zeta(2)$ up to $30$]
The partial sum $\sum_{n\le 30}n^{-2}$ is about $1.612$. The product $\prod_{p\le 30}(1-p^{-2})^{-1}$ is about $1.633$. Both lie below $\pi^2/6\approx 1.645$. The product is larger because it already includes every integer whose prime factors are at most $30$, including those greater than $30$.
\end{examplex}

\begin{examplex}[the series for $\mu$ at $s=2$]
Since $\sum\mu(n)n^{-2}=1/\zeta(2)=6/\pi^2\approx 0.608$. The partial sum through $n=10$ is $1-1/4-1/9+0-1/25-1/36+0+0+0+1/100$, which is already near $0.64$. The omitted tail is small because of $n^{-2}$, which is why $s=2$ is a kinder teacher than $s=1$.
\end{examplex}

\section{Legendre's formula for the sifted count}

% t1-pass2 motive fca497a3
Suppose we want to count the positive integers at most $x$ that are divisible by none of the primes up to $z$. Begin with all integers up to $x$. Subtract those divisible by one small prime. Put back those removed twice, subtract those removed three times, and so on. This is inclusion--exclusion, with the Möbius function carrying all the signs. The exact formula is elementary; estimating it efficiently when there are many small primes is the hard part that leads to sieve theory.

Let $z\ge2$ and let $P(z)=\prod_{p\le z}p$. Write $\Phi(x,z)$ for the number of integers $n$ with $1\le n\le x$ and $\gcd\bigl(n,P(z)\bigr)=1$. These are $1$ and the integers up to $x$ that have no prime factor $\le z$.

\begin{theorem}[Legendre]\label{thm:t1-sieve}
For $x\ge0$ and $z\ge2$,
\[
\Phi(x,z)=\sum_{d\mid P(z)}\mu(d)\Bigl\lfloor\frac xd\Bigr\rfloor.
\]
If $x\ge1$ and $z\ge\sqrt x$, then $\Phi(x,z)=1+\pi(x)-\pi(z)$.
\end{theorem}
\begin{proof}
Theorem~\ref{thm:mobius} says that $\sum_{d\mid\gcd(n,P)}\mu(d)$ equals $1$ if $\gcd(n,P)=1$ and equals $0$ otherwise. Summing over $1\le n\le x$,
\[
\Phi(x,z)=\sum_{n\le x}\sum_{d\mid\gcd(n,P)}\mu(d)
=\sum_{d\mid P}\mu(d)\,\#\{n\le x:d\mid n\}
=\sum_{d\mid P}\mu(d)\Bigl\lfloor\frac xd\Bigr\rfloor.
\]
For the second claim, an integer $n$ with $1<n\le x$ and no prime factor $\le z$ cannot be composite: a composite number has a prime factor at most $\sqrt n\le\sqrt x\le z$. The surviving integers are therefore $1$ and the primes in the interval $(z,x]$.
\end{proof}

% t1-pass2 why
\paragraph{Why this works.}
The Möbius sum $\sum_{d\mid\gcd(n,P)}\mu(d)$ equals $1$ when $n$ is coprime to $P$ and $0$ otherwise. That is a perfect indicator for an integer surviving a sieve. Add it over $n\le x$, exchange the finite sums, and each divisor $d$ contributes exactly $\lfloor x/d\rfloor$ multiples. If $z\ge\sqrt x$, every composite at most $x$ has a prime factor no larger than $z$. The survivors are therefore $1$ and the primes above the cutoff.

The identity is exact. It is also, for large $z$, an impractical way to compute $\pi(x)$, because $P(z)$ has $2^{\pi(z)}$ positive divisors and every one of them contributes a term. At $x=100$ and $z=10$ there are $2^4=16$ terms and the arithmetic is pleasant. At $x=10^4$ and $z=100$ there are $2^{25}$ terms. The sieves of Volumes~6 and~14 exist in order to estimate this sum without expanding it.

\begin{figure}[ht]
\centering
\fig{ch14-sieve}
\caption{The integers from $1$ to $30$, sifted by the primes $2$, $3$ and $5$. Eight numbers remain, in agreement with $\Phi(30,5)=8$. The ochre dots mark the multiples of $5$; many of them were already excluded by $2$ or $3$, which is why an inclusion--exclusion must put them back before removing them again.}
\label{fig:t1-sieve30}
\end{figure}

\begin{examplex}[from thirty, and then to one hundred]
The divisors of $P(5)=30$ and the terms $\mu(d)\lfloor30/d\rfloor$ are
\[
\begin{aligned}
&30,-15,-10,-6,+5,+3,+2,-1.
\end{aligned}
\]
Their sum is $8$. The survivors are $1,7,11,13,17,19,23,29$.

For $\pi(100)$, take $z=10$, so $P=210$. Sixteen terms give $\Phi(100,10)=22$. Theorem~\ref{thm:t1-sieve} converts the count into primes: $\pi(10)=4$, hence $\pi(100)=22-1+4=25$. The same arithmetic with a larger primorial is legitimate and rapidly becomes the wrong tool. The identity's job is structural: every sifted count is a M\"obius average of floors.
\end{examplex}

% t1-pass2 examples fca497a3
\begin{examplex}[sifting $30$ by the primes up to $5$]
The divisors of $30$ contribute $30-15-10-6+5+3+2-1=8$. The survivors are $1,7,11,13,17,19,23,29$. Since $\sqrt{30}>5$, the count is not yet $1+\pi(30)-\pi(5)$ by the sufficient condition of the theorem, although for this particular bound the two numbers happen to agree. The mismatch appears as soon as $49$ enters, as the text computes for $\Phi(50,5)$.
\end{examplex}

\begin{examplex}[the floor identity at $x=10$]
One has $\sum_{n\le 10}\mu(n)\lfloor 10/n\rfloor=10-5-3-2+1-1+1=1$. Every integer $m\le 10$ contributes $\sum_{d\mid m}\mu(d)$, which is $1$ only for $m=1$. This is the sieve with $z$ large enough that $P(z)$ is divisible by every integer up to $x$, written without naming the primes.
\end{examplex}

\section{One series at the boundary}

% t1-pass2 motive 1b6f2fb9
The alternating character $\chi_4$ takes values $1,0,-1,0$ in each block of four integers. Its partial sums are bounded, so the series $\sum_{n\ge1}\chi_4(n)/n$ converges by a summation argument. This does not mean we may automatically evaluate the Euler product at $s=1$. The product includes infinitely many primes, and absolute convergence has disappeared. That missing justification is closely connected with Dirichlet's theorem. We end by marking precisely what we can prove here and what analysis must still supply.

The character $\chi_4$ separates the two odd classes. Its series at the edge $s=1$ can be summed by hand, and the sum is visibly positive. That positivity is the model, in the smallest modulus, of the non-vanishing which Dirichlet's theorem requires.

\begin{proposition}\label{prop:t1-chi4series}
The series $\sum_{n\ge1}\chi_4(n)/n$ converges, and its sum lies in the interval $\bigl(\frac23,1\bigr)$.
\end{proposition}
\begin{proof}
The even terms vanish, and the series is the grouped sum
\[
\sum_{k\ge0}\Bigl(\frac1{4k+1}-\frac1{4k+3}\Bigr)
=\sum_{k\ge0}\frac2{(4k+1)(4k+3)}.
\]
Every grouped term is positive, so the partial sums $S_N$ of the first $N$ groups increase. Writing them out,
\[
S_N=1-\Bigl(\frac13-\frac15\Bigr)-\cdots-\frac1{4N-1}
\]
when $N\ge2$, and $S_1=1-\frac13$. In every case $S_N<1$. Also $S_N\ge S_1=\frac23$, and the inequality is strict for the limit because the second grouped term is positive. An increasing sequence bounded above converges, and the sum lies in $\bigl(\frac23,1\bigr)$.
\end{proof}

\begin{proposition}\label{prop:t1-charbound}
Let $\chi$ be a non-principal character modulo $m$. For every real $x\ge0$,
\[
\Bigl|\sum_{1\le n\le x}\chi(n)\Bigr|\le\varphi(m).
\]
\end{proposition}
\begin{proof}
Let $N=\lfloor x\rfloor$ and write $N=qm+r$ with $0\le r<m$. Each complete block of $m$ consecutive integers contributes $\sum_{n=1}^m\chi(n)=0$, by Proposition~\ref{prop:t1-charsum}. The remainder is a sum of at most $r$ terms of absolute value at most $1$, and at most $\varphi(m)$ of those terms are non-zero. Hence the absolute value is at most $\varphi(m)$.
\end{proof}

\begin{figure}[ht]
\centering
\fig{ch14-race}
\caption{The number of primes $\le x$ in the two odd classes modulo $4$, for $x$ from $10$ to $100$ in steps of $10$. Squares and the solid path count the class $3$; circles and the dashed path count the class $1$. Up to $100$ the class $3$ stays ahead, $13$ primes against $11$. Figure~\ref{fig:mod4-line} shows the same primes one at a time.}
\label{fig:t1-race}
\end{figure}

\begin{think}
Proposition~\ref{prop:t1-chi4series} shows that a signed sum over all integers converges to a positive number. The Euler product expresses the same function for exponents greater than $1$. What extra fact would turn the positivity into $\sum_{p\equiv3\pmod4}1/p=\infty$ and $\sum_{p\equiv1\pmod4}1/p=\infty$? Try to say which step is missing, before reading on.
\end{think}

The missing step is a passage to the limit. For $s>1$,
\[
\log L(s,\chi_4)=\sum_p-\log\bigl(1-\chi_4(p)p^{-s}\bigr)
=\sum_p\frac{\chi_4(p)}{p^s}+O\Bigl(\sum_p\frac1{p^2}\Bigr),
\]
and the error series converges. If $L(s,\chi_4)$ stayed bounded away from $0$ and infinity as $s\downarrow1$, the prime sum $\sum\chi_4(p)/p^s$ would stay bounded, and an elementary argument would promote that bound into the convergence of $\sum\chi_4(p)/p$. Combined with Euler's divergence $\sum_p1/p=\infty$ from Theorem~\ref{thm:euler}, both residue classes would then diverge. Proposition~\ref{prop:t1-chi4series} computes the series of $\chi_4(n)/n$ over all $n$, not the Euler product at $s=1$. The composites and the tail are still in the sum. Removing them is real work. Volume~6 does it for every character, and the non-vanishing of $L(1,\chi)$ is the step that carries the proof.

Exercise~\ref{ex:mod4} already gives infinitely many primes $\equiv3\pmod4$, by a Euclidean construction, and Exercise~\ref{ex:cyclotomic} gives infinitely many primes $\equiv1\pmod4$. Infinitely many is weaker than a divergent reciprocal sum, and both are weaker than Dirichlet's theorem for a general progression. The character makes the gap visible: one can see exactly which identity is still missing.

\begin{beauty}
Where this leads. Quadratic reciprocity says that the Legendre symbol, a character modulo $p$, is essentially a character modulo the numerator. Dirichlet's theorem says that every character except the principal one leaves $L(1,\chi)$ finite and non-zero, and that this single fact forces primes into every eligible residue class. Class field theory, in Volume~13, is the statement that the characters which govern prime splitting in an abelian extension are characters of this same kind. The sieve identity of Theorem~\ref{thm:t1-sieve} is the inclusion--exclusion which Volumes~6 and~14 approximate when the exact sum has too many terms to open. Three later volumes begin where this product, this sign and this floor-sum leave off.
\end{beauty}

% t1-pass2 examples 1b6f2fb9
\begin{examplex}[three grouped terms]
The grouped series begins
\[
\Bigl(1-\frac13\Bigr)+\Bigl(\frac15-\frac17\Bigr)+\Bigl(\frac19-\frac1{11}\Bigr)
=\frac23+\frac2{35}+\frac2{99}.
\]
Numerically this is about $0.6667+0.0571+0.0202=0.744$. The sum is increasing and still less than $1$, and it has already moved well above $2/3$.
\end{examplex}

\begin{examplex}[a partial character sum]
For $\chi_4$, the partial sums of $\chi_4(1)+\cdots+\chi_4(N)$ are $1,1,0,0,1,1,0,0,\ldots$ and never exceed $1$ in absolute value. Proposition~\ref{prop:t1-charbound} only promised the coarser bound $\varphi(4)=2$. The sharper oscillation is special to modulus $4$; the proof for a general modulus stops at $\varphi(m)$.
\end{examplex}

\subsection*{Proof workshop: a small character as a detector}

Let $\chi_4(n)$ be $0$ for even $n$, $1$ for
$n\equiv1\pmod4$, and $-1$ for
$n\equiv3\pmod4$. Then
\[
\sum_{n=1}^{8}\chi_4(n)
=1+0-1+0+1+0-1+0=0.
\]
Over any complete block of four consecutive
integers, the same cancellation occurs:
there is one class with value $1$, one with
value $-1$, and two with value $0$.

This is a finite identity. Do not replace it
silently with a claim about an infinite sum.
For an Euler product, we first work in a region
such as $s>1$ where absolute convergence permits
rearrangement. We then use complete
multiplicativity to collect the terms by prime
powers. The proof separates the arithmetic
fact about $\chi_4$ from the analytic permission
to reorder a series. Both are necessary.


\section{Exercises}

\begin{exercise}[$\star$]\label{ex:t1-char5}
Make the character table of the units modulo $5$, taking $2$ as primitive root, and identify the Legendre symbol in the table. Check the row orthogonality $\sum_{n=1}^5\chi(n)\overline{\psi(n)}$ for two distinct rows.
\end{exercise}

\begin{exercise}[$\star$]\label{ex:t1-chi8check}
Compute $\bigl(\frac2p\bigr)$ and $\chi_8(p)$ for the odd primes $p\le19$. Then prove $\bigl(\frac2p\bigr)=\chi_8(p)$ for every odd prime, using Corollary~\ref{cor:two} and checking the four classes modulo $8$.
\end{exercise}

\begin{exercise}[$\star\star$]\label{ex:t1-sieve30}
Compute $\Phi(30,5)$ using Theorem~\ref{thm:t1-sieve} and list its eight survivors. Compare it with $1+\pi(30)-\pi(5)$. Why do the two counts agree even though $5<\sqrt{30}$, so the theorem's sufficient hypothesis does not hold? Now compute $\Phi(50,5)$ and compare it with $1+\pi(50)-\pi(5)$. Which composite number causes the difference?
\end{exercise}

\begin{exercise}[$\star\star$]\label{ex:t1-phi-ratio}
Prove that $\displaystyle\sum_{d\mid n}\frac{\mu(d)}d=\frac{\varphi(n)}n$ for every $n\ge1$. Deduce $\varphi(n)=n\prod_{p\mid n}\bigl(1-\frac1p\bigr)$.
\end{exercise}

\begin{exercise}[$\star\star$]\label{ex:t1-mu-floor}
Prove that $\displaystyle\sum_{n\le x}\mu(n)\Bigl\lfloor\frac xn\Bigr\rfloor=1$ for every real $x\ge1$. Check the identity at $x=10$.
\end{exercise}

\begin{exercise}[$\star\star$]\label{ex:t1-squarefree}
Show that $\sum_{d^2\mid n}\mu(d)$ equals $1$ if $n$ is squarefree and equals $0$ otherwise. Use it to write the number of squarefree integers up to $x$ as $\sum_{d\le\sqrt x}\mu(d)\bigl\lfloor x/d^2\bigr\rfloor$.
\end{exercise}

\begin{exercise}[$\star\star\star$]\label{ex:t1-char-series}
Let $\chi$ be a non-principal character modulo $m$, and let $A(N)=\sum_{n=1}^N\chi(n)$.
\begin{enumerate}[label=(\alph*)]
\item Prove Proposition~\ref{prop:t1-charbound}: $|A(N)|\le\varphi(m)$.
\item Prove the summation identity
\[
\sum_{n=1}^N\chi(n)\,b_n=A(N)b_N-\sum_{n=1}^{N-1}A(n)\,(b_{n+1}-b_n)
\]
for any numbers $b_n$.
\item Take $b_n=1/n$ and deduce that $\sum_{n\ge1}\chi(n)/n$ converges.
\end{enumerate}
The series of Proposition~\ref{prop:t1-chi4series} is the case $\chi=\chi_4$. The argument here gives convergence for every non-principal character, and it does not by itself evaluate the sum.
\end{exercise}

\begin{exercise}[$\star$]\label{ex:t1-group3}
Compute the sum of the first three grouped terms of $\sum\chi_4(n)/n$, and check that it lies strictly between $2/3$ and $1$.
\end{exercise}

\begin{exercise}[$\star\star$]\label{pb:pseudo}
An odd composite $n$ with $2^{n-1}\equiv1\pmod n$ is a \emph{pseudoprime} (to base $2$).
\begin{enumerate}[label=(\alph*)]
\item Show that $341=11\cdot31$ is a pseudoprime.
\item Show that if $n$ is a pseudoprime, so is $2^n-1$. Deduce that there are infinitely many pseudoprimes.
\item Show that $\frac{4^p-1}3$ is a pseudoprime for every prime $p\ge5$.
\end{enumerate}
\end{exercise}

\begin{exercise}[$\star\star$]\label{pb:pocklington}
\begin{enumerate}[label=(\alph*)]
\item Let $n>1$, and let $F\mid n-1$. Suppose that $a^{n-1}\equiv1\pmod n$ and that $\gcd(a^{(n-1)/q}-1,n)=1$ for every prime
$q\mid F$. Show that every prime factor of $n$ is $\equiv1\pmod F$. Deduce that $n$ is prime if $F>\sqrt n$.
\item (Proth) Let $N=k2^m+1$ with $k$ odd and $k<2^m$. Show that if $a^{(N-1)/2}\equiv-1\pmod N$ for some $a$, then $N$ is
prime. Check that $3329=13\cdot2^8+1$ is prime with $a=3$.
\end{enumerate}
\end{exercise}

\begin{exercise}[$\star\star$]\label{pb:fermatquot}
For a prime $p$ and $p\nmid a$, the \emph{Fermat quotient} is $q_p(a)=\frac{a^{p-1}-1}p$, an integer.
\begin{enumerate}[label=(\alph*)]
\item Show that $q_p(ab)\equiv q_p(a)+q_p(b)\pmod p$, and $q_p(a+pt)\equiv q_p(a)-\frac ta\pmod p$.
\item Show that $2q_p(2)\equiv\sum_{k=1}^{p-1}\frac{(-1)^{k-1}}k\pmod p$ for odd $p$.
\item (Lerch) Let $W_p=\frac{(p-1)!+1}p$. Show that $\sum_{a=1}^{p-1}q_p(a)\equiv W_p\pmod p$.
\item Show that $2^{p-1}\equiv1\pmod{p^2}$ if and only if $2$ has the same order modulo $p$ and modulo $p^2$.
\end{enumerate}
\end{exercise}

\begin{exercise}[$\star\star$]\label{pb:mullin}
\begin{enumerate}[label=(\alph*)]
\item Let $a_1=2$ and let $a_{n+1}$ be the smallest prime factor of $a_1a_2\cdots a_n+1$. Show that the $a_n$ are distinct, and
compute $a_1,\dots,a_7$.
\item Let $b_1=2$ and let $b_{n+1}$ be the \emph{largest} prime factor of $b_1\cdots b_n+1$. Show that $5$ never appears.
\end{enumerate}
\end{exercise}

\begin{exercise}[$\star$]\label{pb:artin}
Let $q$ be a prime such that $p=2q+1$ is also prime, with $p>7$. Show that $2$ is a primitive root modulo $p$ if
$q\equiv1\pmod4$, and that $-2$ is one if $q\equiv3\pmod4$.
\end{exercise}

\begin{exercise}[$\star\star\star$]\label{pb:fermatcubic}
\emph{Gauss's cubic.} Let $p\equiv1\pmod3$ be prime and $N$ the number of $(x,y)\in\mathbb F_p^2$ with $x^3+y^3=1$.
\begin{enumerate}[label=(\alph*)]
\item Let $\chi$ be a character of order $3$ of $\mathbb F_p^*$, extended by $\chi(0)=0$, and for characters $\alpha,\beta$
let $J(\alpha,\beta)=\sum_{u+v=1}\alpha(u)\beta(v)$ (with the convention that the trivial character takes the value $1$ at $0$).
Show that $N=p-2+2\operatorname{Re}J(\chi,\chi)$.
\item Let $g(\alpha)=\sum_{t\in\mathbb F_p}\alpha(t)e(t)$. Show that $g(\alpha)g(\beta)=J(\alpha,\beta)g(\alpha\beta)$ when
$\alpha\beta$ is not trivial, that $|g(\chi)|^2=p$, and that $g(\chi)^3=p\,J(\chi,\chi)$.
\item Write $\omega=e^{2i\pi/3}$, so that $J(\chi,\chi)=a'+b'\omega$ with $a',b'\in\Z$. Show that $a'\equiv-1$ and
$b'\equiv0\pmod 3$.
\item Conclude: $N=p-2+a$, where $a$ is the integer with $a\equiv1\pmod3$ and $4p=a^2+27b^2$ for some integer $b$. Check $p=7$
and $p=13$, and show that $|N-(p-2)|\le2\sqrt p$.
\end{enumerate}
\end{exercise}

\begin{exercise}[$\star\star$]\label{pb:euclidgame}
\emph{Euclid's game (Cole and Davie, 1969).} Two players start from a pair $(a,b)$ of positive integers. A move replaces the larger
number by itself minus a positive multiple of the smaller, as long as the result is $\ge0$. The player who produces a $0$ wins.
\begin{enumerate}[label=(\alph*)]
\item Let $\varphi=\frac{1+\sqrt5}2$. Show that, for $a<b$, the player to move wins if and only if $\frac ba>\varphi$.
\item Deduce that, for random $a<b$, the first player wins with probability tending to $\frac1\varphi=0.618\ldots$
\end{enumerate}
\end{exercise}

\begin{exercise}[$\star\star$]\label{pb:pascalthree}
\emph{Pascal's triangle modulo $3$.}
\begin{enumerate}[label=(\alph*)]
\item Let $p$ be prime and $n=\sum_in_ip^i$ in base $p$. Show that exactly $\prod_i(n_i+1)$ entries of row $n$ of Pascal's triangle
are not divisible by $p$.
\item For $p=3$, let $a_j$ be the number of digits of $n$ equal to $j$. Show that row $n$ has $2^{a_1}$ more entries $\equiv1$ than
entries $\equiv2\pmod3$.
\item In the first $3^k$ rows, find the numbers of entries $\equiv0$, $\equiv1$ and $\equiv2\pmod3$, and their proportions as
$k\to\infty$.
\end{enumerate}
\end{exercise}

\begin{exercise}[$\star\star$]\label{pb:powerful}
\emph{Powerful numbers.} An integer $n\ge1$ is \emph{powerful} if $p^2\mid n$ for every prime $p\mid n$.
\begin{enumerate}[label=(\alph*)]
\item Show that every powerful $n$ can be written uniquely as $a^2b^3$ with $b$ squarefree.
\item Show that the number of powerful $n\le x$ is $\frac{\zeta(3/2)}{\zeta(3)}\sqrt x+O(x^{1/3})$, and that their Dirichlet series is
$\frac{\zeta(2s)\zeta(3s)}{\zeta(6s)}$.
\item Show that there are infinitely many pairs of consecutive powerful numbers, such as $(8,9)$ and $(288,289)$.
\end{enumerate}
\end{exercise}

\begin{exercise}[Journey]\label{pb:jacobi}
\emph{Jacobi's four-square theorem.} For $|x|<1$ let $\Theta(x)=\sum_{m\in\Z}x^{m^2}$, so that $\Theta(x)^4=\sum_{n\ge0}r_4(n)x^n$, where $r_4(n)$ is the number of
$(a,b,c,d)\in\Z^4$ with $a^2+b^2+c^2+d^2=n$. Write $E(x)=\prod_{n\ge1}(1-x^n)$, and $D=x\frac{\mathrm d}{\mathrm dx}$.
\begin{enumerate}[label=(\alph*)]
\item Show that $\Theta(-x)=\prod_{n\ge1}\frac{1-x^n}{1+x^n}$.
\item Show that $E(x)^6=\frac12\sum_{r,s\in\Z}\bigl((2r+1)^2-4s^2\bigr)x^{r^2+r+s^2}$.
\item Let $B(x)=\sum_{r\in\Z}x^{r^2+r}$. Show that $2E^6=\Theta\,(B+4DB)-4B\,D\Theta$, and that
\[
\Theta(x)=\prod_{n\ge1}(1-x^{2n})(1+x^{2n-1})^2,\qq B(x)=2\prod_{n\ge1}(1-x^{2n})(1+x^{2n})^2.
\]
\item Deduce that
\[
\Theta(-x)^4=1-8\sum_{n\ge1}\Bigl(\frac{(2n-1)x^{2n-1}}{1+x^{2n-1}}-\frac{2nx^{2n}}{1+x^{2n}}\Bigr),
\]
and then that $r_4(n)=8\sum_{d\mid n,\,4\nmid d}d$ for every $n\ge1$.
\item Deduce Lagrange's theorem that every positive integer is a sum of four squares, and compute $r_4(p)$ for an odd prime $p$ and $r_4(2^k)$.
\item By counting lattice points in a ball of $\R^4$, show that $\sum_{n\le x}r_4(n)=\frac{\pi^2}2x^2+O(x^{3/2})$. Show also, from (d), that
$\sum_{n\le x}r_4(n)=3\zeta(2)x^2+O(x\log x)$, and conclude that $\zeta(2)=\frac{\pi^2}6$. (Numerically, one finds $\sum_{n\le10^7}r_4(n)=493\,480\,252\,693\,888$, which is
$1.00000007\cdot\frac{\pi^2}2\,10^{14}$.)
\end{enumerate}
\end{exercise}

\begin{exercise}[Journey]\label{pb:densities}
\emph{Four squares, one prime at a time.} For a prime $p$, $k\ge0$ and $n\ge1$, let $N_{p^k}(n)$ be the number of $x\in(\Z/p^k\Z)^4$ with
$x_1^2+x_2^2+x_3^2+x_4^2\equiv n\pmod{p^k}$ (so $N_1(n)=1$). Call such an $x$ \emph{primitive} if some $x_i$ is not divisible by $p$, and let $N^*_{p^k}(n)$ count the
primitive ones.
\begin{enumerate}[label=(\alph*)]
\item Let $p$ be odd. Show that $N^*_{p^{k+1}}(n)=p^3N^*_{p^k}(n)$ for $k\ge1$, and that for $k\ge2$, $N_{p^k}(n)-N^*_{p^k}(n)$ is $p^4N_{p^{k-2}}(\frac n{p^2})$ if
$p^2\mid n$, and $0$ otherwise.
\item Deduce that the \emph{local density} $\delta_p(n)=\lim_{k\to\infty}p^{-3k}N_{p^k}(n)$ exists, and, using Problem~\ref{pb:sumsquares}, that
$\delta_p(n)=\bigl(1-\frac1{p^2}\bigr)\bigl(1+\frac1p+\dots+\frac1{p^a}\bigr)$ when $p^a$ is the exact power of $p$ dividing $n$.
\item Let $p=2$. Show that $N^*_{2^{k+1}}(n)=8N^*_{2^k}(n)$ for $k\ge3$. (Every lift of a solution modulo $2^k$ has the same sum of squares modulo $2^{k+1}$; pair the
solutions.) Deduce from a count modulo $8$ that $\delta_2(n)=1$ for odd $n$, and $\delta_2(n)=3\cdot2^{-a}$ when $2^a$ exactly divides $n$ and $a\ge1$.
\item Show that for every $n\ge1$,
\[
r_4(n)=\pi^2n\prod_p\delta_p(n),
\]
the product converging absolutely. Interpret the factor $\pi^2n$ as a density at infinity: the derivative of the volume of the ball $x_1^2+\dots+x_4^2\le n$.
\end{enumerate}
\end{exercise}

\begin{figure}[ht]
\centering
\fig{ch05-cubic}
\caption{For every prime $p\equiv1\pmod3$ below $3000$, the error $a=N-(p-2)$ in the number of points of $x^3+y^3=1$, divided by
$2\sqrt p$ (Problem~\ref{pb:fermatcubic}). It never leaves $[-1,1]$, and it does not settle anywhere: the points fill the band
with a density that peaks near the edges.}
\label{fig:cubic}
\end{figure}

\section{Hints}

\begin{hint}{ex:t1-char5}
The powers of $2$ modulo $5$ are $1,2,4,3$. A character is determined by $\chi(2)$, a fourth root of unity. The Legendre symbol is the unique non-principal character with real values.
\end{hint}
\begin{hint}{ex:t1-chi8check}
The four odd classes modulo $8$ are $1,3,5,7$, and $(p^2-1)/8$ is even precisely for $p\equiv\pm1\pmod8$.
\end{hint}
\begin{hint}{ex:t1-sieve30}
The divisors of $30$ are $1,2,3,5,6,10,15,30$. The identity with $\pi$ holds if every composite up to $x$ has a prime factor at most $z$, even if the convenient bound $z\ge\sqrt x$ fails. Compare the smallest possible composite whose prime factors all exceed $5$ with $30$ and with $50$.
\end{hint}
\begin{hint}{ex:t1-phi-ratio}
Proposition~\ref{prop:phisum} says $\sum_{d\mid n}\varphi(d)=n$. Invert it.
\end{hint}
\begin{hint}{ex:t1-mu-floor}
Write $\lfloor x/n\rfloor$ as the number of multiples of $n$ up to $x$, exchange the sums, and apply $\mu*\mathbf1=\varepsilon$.
\end{hint}
\begin{hint}{ex:t1-squarefree}
The sum is multiplicative. At a prime power $p^k$, the divisors $d$ with $d^2\mid p^k$ and $\mu(d)\ne0$ are few enough to list.
\end{hint}
\begin{hint}{ex:t1-char-series}
Complete residue blocks sum to zero. For the identity, compare $A(n)-A(n-1)$ with $\chi(n)$, and telescope $A(n)b_n-A(n-1)b_n$. The differences $1/n-1/(n+1)$ are $1/(n(n+1))$.
\end{hint}

\begin{hint}{ex:t1-group3}
Each group is $1/(4k+1)-1/(4k+3)$ for $k=0,1,2$.
\end{hint}

\begin{hint}{pb:pseudo}
(b) $2^n-1$ is composite when $n$ is, and $2^n\equiv1$ modulo $2^n-1$; how does $n$ divide $2^n-2$? (c) $3N=2^{2p}-1$.
\end{hint}

\begin{hint}{pb:pocklington}
(a) For a prime $p\mid n$, the order of $a$ modulo $p$ divides $n-1$ but not $(n-1)/q$. (b) $F=2^m$.
\end{hint}

\begin{hint}{pb:fermatquot}
(b) Expand $2^p=(1+1)^p$; $\frac1p\binom pk\equiv\frac{(-1)^{k-1}}k\pmod p$. (c) Apply (a) to $(p-1)!$.
\end{hint}

\begin{hint}{pb:mullin}
(b) What would $b_1\cdots b_n+1$ be, if its largest prime factor were $5$? Look at it modulo $4$.
\end{hint}

\begin{hint}{pb:artin}
The order of $2$ divides $2q$; Euler's criterion and the value of $\bigl(\frac2p\bigr)$ decide whether it is $q$.
\end{hint}

\begin{hint}{pb:fermatcubic}
(a) The number of cube roots of $u$ is $1+\chi(u)+\chi^2(u)$. (c) Raise $g(\chi)$ to the cube modulo $3$: in characteristic $3$,
cubing is additive.
\end{hint}

\begin{hint}{pb:euclidgame}
(a) Induction on $a+b$. When $b>2a$ the mover can choose between two positions, and exactly one of them has ratio below
$\varphi$.
\end{hint}

\begin{hint}{pb:pascalthree}
(b) Lucas, and the sign $\varepsilon(1)=1$, $\varepsilon(2)=-1$, which is multiplicative on $(\Z/3\Z)^*$.
\end{hint}

\begin{hint}{pb:powerful}
(a) Even exponents go into $a^2$. (c) Pell's equation $x^2-8y^2=1$.
\end{hint}

\begin{hint}{pb:jacobi}
(a) Put $z=-1$ in the triple product and use $\prod(1+x^n)\prod(1-x^{2n-1})=1$. (b) Multiply Jacobi's identity for $E^3$ by itself and pair $k$ with $l$ through $r=\frac{k+l}2$
or $\frac{k+l-1}2$. (c) Write $E^6$ as a product of a theta series in $x^2$ and a series in $x$, using the triple product with $x\mapsto x^2$ twice; or verify the identity
by differentiating the logarithms. (d) Logarithmic derivatives of the products in (a) and (c). (f) Each point of $\Z^4$ in the ball owns the unit cube in front of it.
\end{hint}

\begin{hint}{pb:densities}
(a) Hensel: a primitive solution has some $x_i$ prime to $p$, and $x_i\mapsto x_i+p^kt$ changes the sum by $2x_ip^kt$. A non-primitive solution is $px'$. (b) Use $N_p(n)$
from Problem~\ref{pb:sumsquares}, and induct on $a$. (c) Odd squares are $1$ modulo $8$. (d) Compare the product with $8\sum_{d\mid n,4\nmid d}d$ prime by prime.
\end{hint}

\section{Solutions}

\begin{solution}{ex:t1-char5}
\idea The unit group modulo $5$ is cyclic of order $4$. Fix a generator and remember that a character is determined by the value it assigns to that generator. The value must be a fourth root of unity. Write the four possibilities, then extend multiplicatively to every unit class. To identify the quadratic character, compare its values with the Legendre symbol on the same classes.
\writeup The powers $2^0,2^1,2^2,2^3$ are $1,2,4,3$. Set $\chi_j(2^k)=i^{jk}$. On the ordered units $1,2,3,4$ the four rows are $(1,1,1,1)$, $(1,i,-i,-1)$, $(1,-1,-1,1)$ and $(1,-i,i,-1)$. Since $5\equiv5\pmod8$, $\bigl(\frac25\bigr)=-1$, and reciprocity gives $\bigl(\frac35\bigr)=\bigl(\frac53\bigr)=\bigl(\frac23\bigr)=-1$. Thus $\bigl(\frac n5\bigr)$ agrees with the real row $\chi_2$. The sum of $\chi_0\overline{\chi_2}$ on the units is $1-1-1+1=0$. For $\chi_1$ and $\chi_2$, using that $\chi_2$ is real,
\[
1\cdot1+i\cdot(-1)+(-i)\cdot(-1)+(-1)\cdot1=1-i+i-1=0.
\]
The third term uses $\chi_1(3)=-i$. Theorem~\ref{thm:t1-ortho} predicts the cancellation; the four-term sum is the check.
\end{solution}

\begin{solution}{ex:t1-chi8check}
\idea Both sides of the identity depend only on the odd prime's residue modulo $8$. There are four such classes, so a four-row table is not merely evidence: it covers every possible case. Compute the character values for each class and compare them with the supplementary law for $2$. Make clear why that finite check is exhaustive; otherwise it looks like a numerical experiment.
\writeup For $p=3,5,7,11,13,17,19$ the classes modulo $8$ are $3,5,7,3,5,1,3$. Corollary~\ref{cor:two} gives $\bigl(\frac2p\bigr)=-1,-1,+1,-1,-1,+1,-1$, and $\chi_8$ produces the same list because it equals $-1$ on the classes $3$ and $5$ and equals $+1$ on the classes $1$ and $7$. In general $(p^2-1)/8$ is even for $p\equiv1,7\pmod8$ and odd for $p\equiv3,5\pmod8$, which is the definition of $\chi_8(p)$.
\end{solution}

\begin{solution}{ex:t1-sieve30}
\idea Write the sifted count as the Möbius sum over the divisors of the product of the small primes. For $30=2\cdot3\cdot5$, there are exactly eight squarefree divisors, so the entire inclusion--exclusion calculation fits on a page. Work it out once without symbols. Then test the hypothesis involving $z$ and $\sqrt x$ on the two given examples, keeping strict and non-strict inequalities distinct.
\writeup The terms $\mu(d)\lfloor30/d\rfloor$ are $30,-15,-10,-6,5,3,2,-1$, and they add to $8$. The survivors are $1,7,11,13,17,19,23,29$. The second claim of the theorem asks for $z\ge\sqrt x$. Here $\sqrt{30}>5$, so that hypothesis fails, and the proof given for the translation into $\pi(x)$ does not apply. The two counts nevertheless agree: $1+\pi(30)-\pi(5)=1+10-3=8$. Every composite $n\le30$ still has a prime factor at most $5$, because the next prime is $7$ and $7^2=49>30$. The convenient sufficient condition $z\ge\sqrt x$ is slightly stronger than the condition actually used, namely that every composite up to $x$ has a prime factor $\le z$. A genuine mismatch appears as soon as a square of a larger prime fits under $x$. For example $\Phi(50,5)=14$, counting $1$ and the integers up to $50$ free of the primes $2,3,5$, including the composite $49$. On the other side $1+\pi(50)-\pi(5)=1+15-3=13$.
\end{solution}

\begin{solution}{ex:t1-phi-ratio}
\idea The identity $\sum_{d\mid n}\varphi(d)=n$ is a divisor sum. In the convolution notation, it says $\mathbf1*\varphi=\operatorname{id}$. You already know which function undoes convolution by $\mathbf1$: the Möbius function. Apply it to both sides, write out the resulting divisor sum and verify a small case. This is the power of introducing an algebraic operation: inversion becomes an ordinary calculation.
\writeup Proposition~\ref{prop:phisum} says $\sum_{d\mid n}\varphi(d)=n$, that is $\varphi*\mathbf1=\operatorname{id}$. Convolution with $\mu$ and Theorem~\ref{thm:mobius} give $\varphi=\mu*\operatorname{id}$, so
\[
\varphi(n)=\sum_{d\mid n}\mu(d)\,\frac nd=n\sum_{d\mid n}\frac{\mu(d)}d.
\]
Dividing by $n$ is the first claim. The sum is multiplicative in $n$. At $n=\prod p^a$ it equals $\prod(1-1/p)$, because on a prime power the only non-zero terms are $d=1$ and $d=p$.
\end{solution}

\begin{solution}{ex:t1-mu-floor}
\idea A floor $\lfloor x/d\rfloor$ counts multiples of $d$ at most $x$. Replace the floor by that count and exchange the two \emph{finite} sums. For each integer $m$ that remains, the inner sum runs over its divisors and equals $\sum_{d\mid m}\mu(d)$. That is $1$ only when $m=1$ and $0$ otherwise. Check the range of $x$ to see whether $1$ is included.
\writeup For $x\ge1$,
\[
\sum_{n\le x}\mu(n)\Bigl\lfloor\frac xn\Bigr\rfloor
=\sum_{n\le x}\mu(n)\sum_{m\le x/n}1
=\sum_{nm\le x}\mu(n)
=\sum_{N\le x}\sum_{n\mid N}\mu(n)
=\sum_{N\le x}\varepsilon(N)=1.
\]
At $x=10$ the non-zero contributions are
\[
10-5-3-2+1-1+1=1.
\]
\end{solution}

\begin{solution}{ex:t1-squarefree}
\idea The condition ``$n$ is squarefree'' is a condition at each prime. A multiplicative sum that vanishes on $p^2$ and equals $1$ on $p$ is the right detector.
\writeup Let $f(n)=\sum_{d^2\mid n}\mu(d)$. If $\gcd(a,b)=1$, the square divisors of $ab$ are products of square divisors, so $f$ is multiplicative. At $p^k$, the integers $d$ with $d^2\mid p^k$ and $\mu(d)\ne0$ are $d=1$, and also $d=p$ when $k\ge2$. Thus $f(p^k)=1$ for $k=0,1$ and $f(p^k)=1+\mu(p)=0$ for $k\ge2$. Therefore $f(n)=1$ if every exponent in $n$ is $0$ or $1$, and $f(n)=0$ otherwise. Counting,
\begin{align*}
\sum_{n\le x}f(n)
&=\sum_{n\le x}\sum_{d^2\mid n}\mu(d)\\
&=\sum_{d\le\sqrt x}\mu(d)\Bigl\lfloor\frac x{d^2}\Bigr\rfloor,
\end{align*}
which is the number of squarefree positive integers up to $x$.
\end{solution}

\begin{solution}{ex:t1-char-series}
\idea Complete periods cancel. Summation by parts then trades the oscillating coefficients $\chi(n)$ for the bounded sums $A(n)$, at the price of a difference of the sequence $b_n$.
\writeup (a) This is the proof of Proposition~\ref{prop:t1-charbound}.

(b) Since $\chi(n)=A(n)-A(n-1)$ with $A(0)=0$,
\begin{align*}
\sum_{n=1}^N\chi(n)b_n
&=\sum_{n=1}^N\bigl(A(n)-A(n-1)\bigr)b_n\\
&=\sum_{n=1}^N A(n)b_n-\sum_{n=1}^{N-1}A(n)b_{n+1}\\
&=A(N)b_N-\sum_{n=1}^{N-1}A(n)(b_{n+1}-b_n).
\end{align*}

(c) With $b_n=1/n$, one has $b_{n+1}-b_n=-1/(n(n+1))$, so
\[
\sum_{n=1}^N\frac{\chi(n)}n=\frac{A(N)}N+\sum_{n=1}^{N-1}\frac{A(n)}{n(n+1)}.
\]
The first term tends to $0$ because $A(N)$ is bounded. The series $\sum A(n)/(n(n+1))$ converges absolutely, by comparison with $\varphi(m)\sum n^{-2}$. Therefore the partial sums on the left converge.
\end{solution}

\begin{solution}{ex:t1-group3}
\idea The grouped terms are positive, the first is $2/3$, and the partial sum of an increasing sequence bounded by $1$ stays in between once a second positive term is added.
\writeup The three terms are $1-1/3=2/3$, $1/5-1/7=2/35$ and $1/9-1/11=2/99$. Their sum is $2/3+2/35+2/99$. The common denominator is $3465$, and the sum is $(2310+198+70)/3465=2578/3465$. Since $2578/3465-2/3=(7734-6930)/10395>0$ and $1-2578/3465>0$, the sum lies strictly between $2/3$ and $1$.
\end{solution}

\begin{solution}{pb:pseudo}
\idea For (b), the number $m=2^n-1$ satisfies $2^n\equiv1\pmod m$, so $2^{m-1}\equiv1$ as soon as $n\mid m-1$. And
$m-1=2(2^{n-1}-1)$, which $n$ divides exactly because $n$ is a pseudoprime. For (c), the same mechanism with $2^{2p}$.

\writeup (a) $2^{10}=1024=93\cdot11+1$ and $2^5=32=31+1$, so $2^{10}\equiv1\pmod{11}$ and $2^5\equiv1\pmod{31}$. Since $10\mid340$ and
$5\mid340$, $2^{340}\equiv1$ modulo $11$ and $31$, hence modulo $341$.

(b) Let $n$ be a pseudoprime and $m=2^n-1$, which is odd. If $n=ab$ with $1<a<n$, then $2^a-1$ divides $2^n-1$ and lies strictly
between $1$ and $m$, so $m$ is composite. Now $m-1=2(2^{n-1}-1)$ and $n\mid2^{n-1}-1$, so $m-1=nk$ for an integer $k$. As
$2^n\equiv1\pmod m$, $2^{m-1}=(2^n)^k\equiv1\pmod m$. So $m$ is a pseudoprime, larger than $n$. Starting from $341$, this gives
infinitely many.

(c) Let $N=\frac{4^p-1}3=(2^p-1)\cdot\frac{2^p+1}3$. Both factors are integers
larger than $1$ ($3\mid2^p+1$ as $p$ is odd, and $\frac{2^p+1}3>1$ for $p\ge3$). So $N$ is composite, and odd. Next,
$N-1=\frac{4^p-4}3=4\cdot\frac{4^{p-1}-1}3$. By Fermat, $p\mid4^{p-1}-1$, and as $p\ne3$, $p$ divides $\frac{4^{p-1}-1}3$. So $2p\mid N-1$.
Finally $3N=2^{2p}-1$, so $2^{2p}\equiv1\pmod N$ and $2^{N-1}\equiv1\pmod N$.
\tools Volume~2, Chapters~1 and~3 (divisibility, Fermat's little theorem).
\leads (b) is due to Malo (1903) and (c) to Cipolla (1904). Pseudoprimes are rare: Pomerance proved that their number up to $x$
is at most $x\,e^{-\frac12\log x\log\log\log x/\log\log x}$ for large $x$, so that a random odd number that passes the base $2$
test is almost surely prime, which is what software relied on for decades.
\end{solution}

\begin{solution}{pb:pocklington}
\idea Lucas's criterion asks to factor all of $n-1$. Pocklington noticed that a partial factorisation is enough: every prime
factor of $n$ is then $\equiv1$ modulo the factored part $F$, hence larger than $F$, and a composite $n$ has a prime factor at
most $\sqrt n$. For Proth numbers $N-1=k2^m$, the power of $2$ is the factored part.

\writeup (a) Let $p$ be a prime factor of $n$ and $e$ the order of $a$ modulo $p$ ($p\nmid a$ since $a^{n-1}\equiv1$). Then
$e\mid n-1$. Let $q^f$ be the exact power of a prime $q\mid F$ in $F$, and $q^g$ its power in $n-1$, $g\ge f$. If $q^g\nmid e$, then
$e\mid\frac{n-1}q$ and $p\mid a^{(n-1)/q}-1$, contradicting $\gcd(a^{(n-1)/q}-1,n)=1$. So $q^g\mid e\mid p-1$. Hence $F\mid p-1$ and
$p>F$. If $F>\sqrt n$, every prime factor of $n$ exceeds $\sqrt n$, so $n$ is prime.

(b) Take $F=2^m$, with only $q=2$. From $a^{(N-1)/2}\equiv-1$: $a^{N-1}\equiv1$, and $\gcd(a^{(N-1)/2}-1,N)=\gcd(-2,N)=1$ as $N$ is odd. And
$N=k2^m+1\le(2^m-1)2^m+1<2^{2m}$, so $F>\sqrt N$. By (a), $N$ is prime. For $3329$: $m=8$, $k=13<256$, and repeated squaring
gives $3^{1664}\equiv3328\equiv-1\pmod{3329}$, so $3329$ is prime.
\tools Volume~2, Chapters~3 and~4 (fast exponentiation and orders).
\leads Pocklington (1914) and Proth (1878). The largest known primes that are not Mersenne primes are Proth primes, found by the
PrimeGrid project with exactly this test. Note the asymmetry: (b) proves primality in one exponentiation, but a failure says
nothing unless $\bigl(\frac aN\bigr)=-1$.
\end{solution}

\begin{solution}{pb:fermatquot}
\idea $a^{p-1}=1+pq_p(a)$: the Fermat quotient is the first digit of $a^{p-1}$ in base $p$ beyond the units digit, and
multiplying numbers of the form $1+px$ adds the $x$'s modulo $p$. So $q_p$ behaves like a logarithm. For (b), expand $2^p$ by the
binomial theorem; for (c), take the logarithm of Wilson's theorem.

\writeup All congruences are modulo $p$ unless stated. (a) $(ab)^{p-1}=(1+pq_p(a))(1+pq_p(b))\equiv1+p(q_p(a)+q_p(b))\pmod{p^2}$. Next,
$(a+pt)^{p-1}\equiv a^{p-1}+(p-1)pta^{p-2}\pmod{p^2}$, so $q_p(a+pt)\equiv q_p(a)-ta^{p-2}\equiv q_p(a)-\frac ta$, since $a^{p-1}\equiv1$.

(b) $2^p=\sum_k\binom pk=2+\sum_{k=1}^{p-1}\binom pk$, and $\frac1p\binom pk=\frac{(p-1)(p-2)\cdots(p-k+1)}{k!}\equiv\frac{(-1)^{k-1}(k-1)!}{k!}=\frac{(-1)^{k-1}}k$.
So $\frac{2^p-2}p\equiv\sum_{k=1}^{p-1}\frac{(-1)^{k-1}}k$, and $\frac{2^p-2}p=2q_p(2)$.

(c) Extend $q_p(N)=\frac{N^{p-1}-1}p$ to all $N$ prime to $p$; (a) still holds. Then
$\sum_{a=1}^{p-1}q_p(a)\equiv q_p\bigl((p-1)!\bigr)$. Write $(p-1)!=-1+pW_p$ (Wilson). As $p-1$ is even,
$(-1+pW_p)^{p-1}\equiv1-(p-1)pW_p\equiv1+pW_p\pmod{p^2}$, so $q_p((p-1)!)\equiv W_p$.

(d) Let $e$ and $E$ be the orders of $2$ modulo $p$ and $p^2$. Then $e\mid E$, and $2^{pe}=(1+cp)^p\equiv1\pmod{p^2}$ if $2^e=1+cp$, so
$E\mid pe$: $E\in\{e,pe\}$. If $2^{p-1}\equiv1\pmod{p^2}$, then $E\mid p-1$ is prime to $p$, so $E=e$. Conversely if $E=e$, then
$E\mid p-1$ and $2^{p-1}\equiv1\pmod{p^2}$.
\tools Volume~2, Chapters~3 and~4 (Wilson, Fermat, orders); Volume~1, Chapter~1 (binomial theorem).
\leads A prime with $q_p(2)\equiv0$ is a \emph{Wieferich prime}. Wieferich proved in 1909 that if $x^p+y^p=z^p$ has a solution with
$p\nmid xyz$, then $p$ is one. Only $1093$ (Meissner, 1913) and $3511$ (Beeger, 1922) are known; heuristically there should be about
$\log\log x$ of them below $x$, infinitely many but very few. Lerch's formula (c) connects them with Wilson primes ($W_p\equiv0$):
$5$, $13$ and $563$.
\end{solution}

\begin{solution}{pb:mullin}
\idea Euclid's argument guarantees that each new term is a new prime. For (b), the number $P+1$ with $P=b_1\cdots b_n$ is prime to
$2$ and $3$ as soon as $3$ has appeared; if its largest prime factor were $5$, it would be a power of $5$, and $5^c-1$ is
divisible by $4$, whereas $P$ contains the prime $2$ only once.

\writeup (a) $a_{n+1}$ divides $a_1\cdots a_n+1$, which is prime to each $a_i$, so $a_{n+1}\ne a_i$. The products give:
$2+1=3$; $6+1=7$; $42+1=43$; $1806+1=13\cdot139$, so $a_5=13$; $23478+1=23479=53\cdot443$, so $a_6=53$;
$23478\cdot53+1=1244335=5\cdot248867$, so $a_7=5$.

(b) The terms are distinct for the same reason, so for $n\ge1$, $P_n=b_1\cdots b_n$ contains the factor $2$ exactly once. $b_2=3$. For
$n\ge2$, $6\mid P_n$, so $P_n+1$ is prime to $2$ and $3$. If its largest prime factor were $5$, then $P_n+1=5^c$ with $c\ge1$, and
$P_n=5^c-1\equiv0\pmod4$, contradicting $v_2(P_n)=1$. And $b_2=3\ne5$. So $5$ never appears.
\tools Volume~2, Chapters~1 and~2 (Euclid's argument, factorisation).
\leads Mullin asked in 1963 whether every prime appears in the first sequence; Shanks conjectured yes, and it is open, the
sequence being known for only about fifty terms because each one needs a hard factorisation. For the second sequence, Cox and
van der Poorten (1968) showed that $5,11,13,17,19,23,29,31,37,41,47,53$ never appear, by arguments like (b) with more
congruences, and Booker (2012) proved that it misses infinitely many primes.
\end{solution}

\begin{solution}{pb:artin}
\idea The group $(\Z/p\Z)^*$ has order $2q$, so the order of $2$ is $1$, $2$, $q$ or $2q$. Orders $1$ and $2$ are excluded by size. Order
$q$ means $2^q=2^{(p-1)/2}\equiv1$, which by Euler's criterion means that $2$ is a square; and whether $2$ is a square depends on $p$
modulo $8$, which is fixed by $q$ modulo $4$.

\writeup The order $e$ of $2$ divides $p-1=2q$. $e=1$ or $2$ would need $p\mid1$ or $p\mid3$, excluded. If $q\equiv1\pmod4$, $p=2q+1\equiv3\pmod8$,
so $\bigl(\frac2p\bigr)=-1$ and $2^q\equiv-1$ by Euler's criterion: $e\ne q$, so $e=2q$ and $2$ is a primitive root. If $q\equiv3\pmod4$, then
$p\equiv7\pmod8$ and $\bigl(\frac2p\bigr)=1$, $\bigl(\frac{-1}p\bigr)=-1$, so $\bigl(\frac{-2}p\bigr)=-1$; the same argument applies to $-2$, whose order is
not $1$ or $2$ since $p>3$.
\tools Volume~2, Chapters~4 and~5.
\leads Artin conjectured (1927) that $2$ is a primitive root modulo infinitely many primes, in fact for a proportion $\prod_p\bigl(1-\frac1{p(p-1)}\bigr)
=0.3739\ldots$ of them. Hooley (1967) proved it under the generalised Riemann hypothesis. Unconditionally, Heath-Brown (1986) showed that at most two
primes fail to be primitive roots for infinitely many $p$: so one of $2$, $3$, $5$ works, but nobody knows which. This problem would
give a proof if there were infinitely many such ``safe primes'', which is also open.
\end{solution}

\begin{solution}{pb:fermatcubic}
\idea Counting solutions of $u+v=1$ with $u,v$ cubes is a sum of Jacobi sums. The only non-trivial one is $J(\chi,\chi)$, an
integer of $\Z[\omega]$ of norm $p$, and a congruence modulo $3$ picks it out among the six associates.

\writeup (a) Since $p\equiv1\pmod 3$, $\mathbb F_p^*$ is cyclic of order divisible by $3$; it has two characters $\chi,\chi^2=\bar\chi$
of order $3$. For $u\ne0$, the number of $x$ with $x^3=u$ is $3$ if $u$ is a cube and $0$ otherwise, which is $1+\chi(u)+\chi^2(u)$;
for $u=0$ it is $1$, which is also $1+\chi(0)+\chi^2(0)$ with our conventions. Writing $\varepsilon$ for the trivial character,
\[
N=\sum_{u+v=1}\Bigl(\sum_{i=0}^2\chi^i(u)\Bigr)\Bigl(\sum_{j=0}^2\chi^j(v)\Bigr)=\sum_{i,j}J(\chi^i,\chi^j).
\]
Now $J(\varepsilon,\varepsilon)=p$ and $J(\varepsilon,\psi)=\sum_v\psi(v)=0$ for $\psi$ non-trivial. Next
$J(\chi,\bar\chi)=\sum_{u\ne0,1}\chi\bigl(\frac u{1-u}\bigr)$, and $u\mapsto\frac u{1-u}$ is a bijection from
$\mathbb F_p\setminus\{0,1\}$ to $\mathbb F_p\setminus\{0,-1\}$; so $J(\chi,\bar\chi)=-\chi(-1)=-1$, because
$\chi(-1)^2=\chi(1)=1$ and $\chi(-1)^3=1$ force $\chi(-1)=1$. The same holds for $J(\bar\chi,\chi)$. Finally
$J(\bar\chi,\bar\chi)=\overline{J(\chi,\chi)}$. So $N=p-2+2\operatorname{Re}J(\chi,\chi)$.

(b) Expanding,
\[
g(\alpha)g(\beta)=\sum_s e(s)\sum_{x+y=s}\alpha(x)\beta(y).
\]
For $s\ne0$, put $x=su$, $y=sv$: the inner sum is $\alpha\beta(s)J(\alpha,\beta)$. For $s=0$ it is
\[
\sum_x\alpha(x)\beta(-x)=\beta(-1)\sum_x\alpha\beta(x)=0,
\]
since $\alpha\beta$ is non-trivial. So
\[
g(\alpha)g(\beta)=J(\alpha,\beta)\sum_s\alpha\beta(s)e(s)=J(\alpha,\beta)g(\alpha\beta).
\]

Next, $|g(\chi)|^2=\sum_{x,y\ne0}\chi(x)\bar\chi(y)e(x-y)$. Put $x=ty$: this is
$\sum_{t\ne0}\chi(t)\sum_{y\ne0}e(y(t-1))$. The inner sum is $p-1$ for $t=1$ and $-1$ otherwise, so
$|g|^2=(p-1)-\sum_{t\ne0,1}\chi(t)=(p-1)+1=p$. Also $\overline{g(\chi)}=\sum\bar\chi(t)e(-t)=\chi(-1)g(\bar\chi)=g(\bar\chi)$.
Apply the first identity twice: $g(\chi)^2=J(\chi,\chi)g(\chi^2)=J(\chi,\chi)g(\bar\chi)$, so
$g(\chi)^3=J(\chi,\chi)g(\bar\chi)g(\chi)=J(\chi,\chi)|g(\chi)|^2=pJ(\chi,\chi)$.

(c) Work in the ring $R=\Z[\omega,\zeta]$ with $\zeta=e^{2i\pi/p}$. Modulo $3R$, $(a+b)^3\equiv a^3+b^3$ (the middle binomial
coefficients are $3$), so
\[
g(\chi)^3\equiv\sum_t\chi(t)^3\zeta^{3t}=\sum_{t\ne0}\zeta^{3t}=-1\pmod{3R},
\]
because $\chi(t)^3=1$ for $t\ne0$ and $3t$ runs over $\mathbb F_p^*$. Since $p\equiv1\pmod3$, (b) gives $J(\chi,\chi)\equiv-1$,
that is $\frac{J+1}3\in R$. But $\frac{J+1}3=\frac{a'+1}3+\frac{b'}3\omega$ is in $\Q(\omega)$ and is an algebraic integer, so it lies
in $\Z[\omega]$; as $1,\omega$ is a basis of $\Z[\omega]$ and of $\Q(\omega)$, $3\mid a'+1$ and $3\mid b'$.

(d) Let $a=2\operatorname{Re}J=2a'-b'$ (because $\operatorname{Re}\omega=-\frac12$). Then $a\equiv-2\equiv1\pmod 3$. By (b),
$|J|^2=\frac{|g|^6}{p^2}=p$, and $|a'+b'\omega|^2=a'^2-a'b'+b'^2$, so $4p=(2a'-b')^2+3b'^2=a^2+27b^2$ with $b=\frac{b'}3$. And
$N=p-2+a$. Since $27b^2\ge0$, $|a|\le2\sqrt p$.

For $p=7$: $28=1+27$, $a=1$, $N=6$. Directly: the cubes mod $7$ are $0,1,6$, and $u+v=1$ with $u,v$ cubes forces $\{u,v\}=\{0,1\}$,
giving $1\cdot3+3\cdot1=6$ solutions. For $p=13$: $52=25+27$, $a\equiv1\pmod3$ picks $a=-5$, so $N=6$. (Why is $a$ determined?
If $4p=a^2+27b^2=c^2+27d^2$, the numbers $\frac{a\pm3b\sqrt{-3}}2$ and $\frac{c\pm3d\sqrt{-3}}2$ have norm $p$ in the unique
factorisation ring $\Z[\omega]$, so they are associates, and the congruence modulo $3$ leaves only $c=a$.)

\tools Volume~2, Chapters~4 and~5 (characters, Gauss sums); Volume~2, Chapter~10 ($\Z[\omega]$) and Volume~13, Chapter~1 (algebraic integers); Volume~3, Chapter~9 (Frobenius).
\leads This is Gauss's result in article~358 of the \emph{Disquisitiones} (1801), and André Weil's 1949 paper on equations over
finite fields starts from it: the projective curve $x^3+y^3=z^3$ has $p+1-a$ points, and $|a|\le2\sqrt p$ is the Hasse bound for
this elliptic curve. Figure~\ref{fig:cubic} shows $\frac a{2\sqrt p}$; because the curve has complex multiplication by $\Z[\omega]$,
its distribution is not the Sato--Tate semicircle but the arcsine-like law of $\cos\theta$ for uniform $\theta$, which piles up at
$\pm1$.
\end{solution}

\begin{solution}{pb:euclidgame}
\idea When $b<2a$ there is only one move, so nothing is decided by the mover. When $b>2a$ the mover can stop at
$(a,a+r)$ or go on to $(r,a)$, and $\varphi$ is exactly the ratio that separates the winning choice from the losing one,
because $\varphi-1=\frac1\varphi$.

\writeup (a) Call a position $(a,b)$ with $a\le b$ \emph{winning} if the player to move can force a win. If $a=b$ the mover wins at once.
We show by induction on $a+b$ that for $a<b$, $(a,b)$ is winning if and only if $b/a>\varphi$ (equality is impossible, $\varphi$ being
irrational). Write $b=qa+r$ with $0\le r<a$, $q\ge1$.

\emph{Case $b/a<\varphi$.} Then $q=1$ (as $\varphi<2$) and $r=b-a>0$, the only move gives $(r,a)$, and $\frac ar>\varphi$ because
$\frac ra=\frac ba-1<\varphi-1=\frac1\varphi$. By induction the opponent, now to move from $(r,a)$, wins.

\emph{Case $b/a>\varphi$.} If $r=0$ the mover wins at once by producing a $0$. Otherwise, if $q=1$, the only move
gives $(r,a)$ with $\frac ar<\varphi$ (since $\frac ra>\frac1\varphi$), a losing position for the opponent. If $q\ge2$, the mover may
leave either $(r,a)$ or $(a,a+r)$. Now $\frac{a+r}a<\varphi$ if and only if $\frac ra<\frac1\varphi$, if and only if $\frac ar>\varphi$; so
exactly one of the two positions has ratio below $\varphi$, and the mover picks it. By induction the opponent loses.

(b) The number of pairs $1\le a<b\le N$ with $b>\varphi a$ is $\sum_{a<N/\varphi}(N-\lfloor\varphi a\rfloor)=\frac{N^2}{2\varphi}+O(N)$,
out of $\binom N2$ pairs; the ratio tends to $\frac1\varphi$.

\tools Volume~2, Chapters~1 and~8 (Euclid's algorithm, continued fractions); Volume~1, Chapter~3 (the golden ratio).
\leads The game follows Euclid's algorithm, and the only real decisions are at partial quotients $\ge2$: the winner is decided by
the first partial quotient of $b/a$ that is at least $2$, by its parity. Cole and Davie's paper appeared in the \emph{Mathematical
Gazette} 53 (1969).
\end{solution}

\begin{solution}{pb:pascalthree}
\idea Lucas's theorem makes $\binom nk\bmod p$ a product of digit-by-digit factors, so any multiplicative statistic of the row
factors over the digits.

\writeup (a) By Lucas, $\binom nk\not\equiv0$ if and only if $k_i\le n_i$ for every digit, and there are $\prod_i(n_i+1)$ such $k$.

(b) Let $\varepsilon(1)=1$, $\varepsilon(2)=-1$; it is a homomorphism $(\Z/3\Z)^*\to\{\pm1\}$. Summing over the $k$ of (a),
\[
\#\{k:\tbinom nk\equiv1\}-\#\{k:\tbinom nk\equiv2\}=\sum_{k_i\le n_i}\varepsilon\Bigl(\prod_i\tbinom{n_i}{k_i}\Bigr)=\prod_i\sum_{k_i=0}^{n_i}\varepsilon\tbinom{n_i}{k_i}.
\]
The inner sum is $1$ for $n_i=0$, $1+1=2$ for $n_i=1$, and $1-1+1=1$ for $n_i=2$ (as $\binom21=2$). So the difference is $2^{a_1}$.

(c) Rows $0,\dots,3^k-1$ are the $n$ with $k$ digits $n_i\in\{0,1,2\}$. Summing the products over all choices,
the number of non-zero entries is $(1+2+3)^k=6^k$ and the difference between the $1$'s and the $2$'s is $(1+2+1)^k=4^k$. So there
are $\frac{6^k+4^k}2$ entries $\equiv1$, $\frac{6^k-4^k}2$ entries $\equiv2$, and $\frac{3^k(3^k+1)}2-6^k$ multiples of $3$. The proportion of
multiples of $3$ tends to $1$, while among the others the $1$'s and $2$'s become equally frequent, since $\frac{4^k}{6^k}\to0$.

\tools Exercise~3.9 (Lucas's theorem).
\leads For a general prime $p$, the multiplicative characters of $(\Z/p\Z)^*$ play the role of $\varepsilon$, and their sums over a
row give the number of entries in each residue class. This was worked out by Garfield and Wilf (1992); Granville's survey
listed below tells the story and does the same modulo prime powers.
\end{solution}

\begin{solution}{pb:powerful}
\idea A powerful number has every exponent $\ge2$, and every such exponent is $2u+3v$ with $v\in\{0,1\}$. So powerful numbers are
squares times cubes of squarefree numbers, and counting them is counting lattice points under a curve.

\writeup (a) Let $n=\prod p^{e_p}$ with all $e_p\ge2$. Put $v_p=1$ if $e_p$ is odd and $v_p=0$ otherwise, and $u_p=\frac{e_p-3v_p}2\ge0$. Then
$n=a^2b^3$ with $a=\prod p^{u_p}$ and $b=\prod p^{v_p}$ squarefree. Conversely $n=a^2b^3$ with $b$ squarefree forces $v_p\equiv e_p\pmod2$,
so the decomposition is unique, and every $a^2b^3$ is powerful.

(b) By (a) the count is
\[
\sum_{b\le x^{1/3}}\mu^2(b)\Bigl\lfloor\sqrt{x/b^3}\Bigr\rfloor=\sqrt x\sum_{b\le x^{1/3}}\frac{\mu^2(b)}{b^{3/2}}+O(x^{1/3}).
\]
Since $\sum_{b>y}b^{-3/2}=O(y^{-1/2})$, completing the sum costs $O(\sqrt x\,x^{-1/6})=O(x^{1/3})$, and
$\sum_b\mu^2(b)b^{-3/2}=\frac{\zeta(3/2)}{\zeta(3)}$ because $\sum\mu^2(b)b^{-s}=\prod_p(1+p^{-s})=\frac{\zeta(s)}{\zeta(2s)}$. For the Dirichlet series,
the factor at $p$ is, with $y=p^{-s}$,
\[
1+\sum_{e\ge2}y^e=\frac{1-y+y^2}{1-y},\qquad\text{and}\qquad\frac{1-y^6}{(1-y^2)(1-y^3)}=\frac{1+y^3}{1-y^2}=\frac{(1+y)(1-y+y^2)}{(1-y)(1+y)},
\]
the same.

(c) Pell's equation $X^2-8Y^2=1$ has infinitely many solutions (Theorem~8.8), for instance $(3,1)$ and $(17,6)$, the
others coming from powers of $3+\sqrt8$. For each, $8Y^2=2^3Y^2$ is powerful (every prime of $Y$ appears at least squared, and $2$ at
least cubed) and $8Y^2+1=X^2$ is a square. This gives $(8,9)$ and $(288,289)$.

\tools Theorem~8.8 (Pell's equation); Volume~6, Chapter~2 (counting with Dirichlet series).
\leads Erdős conjectured that there are no three consecutive powerful numbers; this is still open (and would follow from the
$abc$ conjecture). The next term of the asymptotic expansion is $\frac{\zeta(2/3)}{\zeta(2)}x^{1/3}$, negative since $\zeta(2/3)<0$, which is why
the numerical counts lie below $2.173\sqrt x$.
\end{solution}

\begin{solution}{pb:jacobi}
\idea Jacobi's identity squares to a double sum, and a change of variables turns that double sum into a combination of two theta series and their derivatives. Dividing
by the theta series turns the derivatives into logarithmic derivatives of products, which are Lambert series: divisor sums appear by themselves.

\writeup (a) Put $z=-1$ in the triple product: $\Theta(-x)=\sum(-1)^mx^{m^2}=\prod(1-x^{2n})(1-x^{2n-1})^2$. Since $\prod(1-x^{2n-1})=\frac{\prod(1-x^n)}{\prod(1-x^{2n})}=\frac1{\prod(1+x^n)}$
and $1-x^{2n}=(1-x^n)(1+x^n)$, the product is $\prod(1-x^n)(1+x^n)\cdot\prod(1+x^n)^{-2}=\prod\frac{1-x^n}{1+x^n}$.

(b) Replacing $k$ by $-1-k$ leaves $(-1)^k(2k+1)$ and $\frac{k(k+1)}2$ unchanged, so $E^3=\frac12\sum_{k\in\Z}(-1)^k(2k+1)x^{k(k+1)/2}$ and
\[
E^6=\frac14\sum_{k,l\in\Z}(-1)^{k+l}(2k+1)(2l+1)x^{(k^2+k+l^2+l)/2}.
\]
If $k+l$ is even, put $k=r+s$, $l=r-s$: then $(2k+1)(2l+1)=(2r+1)^2-4s^2$ and $\frac{k^2+k+l^2+l}2=r^2+r+s^2$. If $k+l$ is odd, put $k=s+r$, $l=s-r-1$: then the sign is $-1$,
$(2k+1)(2l+1)=4s^2-(2r+1)^2$ and the exponent is again $r^2+r+s^2$. Both substitutions are bijections onto $\Z^2$, so each $(r,s)$ receives $2\bigl((2r+1)^2-4s^2\bigr)$, and
$E^6=\frac12\sum_{r,s}\bigl((2r+1)^2-4s^2\bigr)x^{r^2+r+s^2}$.

(c) Since $(2r+1)^2=4(r^2+r)+1$, the double sum of (b) is $\sum_r\bigl(4(r^2+r)+1\bigr)x^{r^2+r}\cdot\sum_sx^{s^2}-4\sum_rx^{r^2+r}\cdot\sum_ss^2x^{s^2}=(4DB+B)\Theta-4B\,D\Theta$.
The triple product with $z=1$ gives the product for $\Theta$. With $z=x$ it gives $\sum_mx^{m^2+m}=\prod(1-x^{2n})(1+x^{2n})(1+x^{2n-2})$, and the factor $n=1$ of the last product is
$2$, so $B=2\prod(1-x^{2n})(1+x^{2n})^2$.

(d) Divide (c) by $\Theta B$: $\frac{2E^6}{\Theta B}=1+4\frac{DB}B-4\frac{D\Theta}\Theta$. By the products,
$\Theta B=2\prod(1-x^{2n})^2(1+x^n)^2$, so $\frac{2E^6}{\Theta B}=\prod\frac{(1-x^n)^6}{(1-x^n)^2(1+x^n)^4}=\Theta(-x)^4$ by (a). Logarithmic derivatives of the products give
\[
\frac{DB}B-\frac{D\Theta}\Theta=\sum_{n\ge1}\Bigl(\frac{4nx^{2n}}{1+x^{2n}}-\frac{2(2n-1)x^{2n-1}}{1+x^{2n-1}}\Bigr),
\]
the terms $-\frac{2nx^{2n}}{1-x^{2n}}$ cancelling. This is the formula for $\Theta(-x)^4$. Now replace $x$ by $-x$:
\[
\Theta(x)^4=1+8\sum_{m\text{ odd}}\frac{mx^m}{1-x^m}+8\sum_{m\text{ even}}\frac{mx^m}{1+x^m}.
\]
Expanding, $r_4(N)=8\sum_{m\mid N}m\,\varepsilon(m)$ with $\varepsilon=1$ for odd $m$ and $\varepsilon=(-1)^{N/m-1}$ for even $m$. For odd $N$ this is $8\sigma(N)$. For
$N=2^au$, $a\ge1$, $u$ odd, the even divisors $2^bv$ ($v\mid u$) have sign $-1$ for $b<a$ and $+1$ for $b=a$, so the total is
$8\sigma(u)\bigl(1+2^a-(2+4+\dots+2^{a-1})\bigr)=8\cdot3\sigma(u)$, which is $8\sum_{d\mid N,4\nmid d}d$ again.

(e) $r_4(n)\ge8>0$, so every $n\ge1$ is a sum of four squares. $r_4(p)=8(p+1)$, and $r_4(2^k)=8(1+2)=24$ for $k\ge1$.

(f) Let $R=\sqrt x$. The points $v\in\Z^4$ with $|v|\le R$ number $1+\sum_{n\le x}r_4(n)$, and the unit cubes $v+[0,1[^4$ cover the ball of radius $R-2$ and lie in the
ball of radius $R+2$. The volumes are $\frac{\pi^2}2(R\pm2)^4=\frac{\pi^2}2x^2+O(x^{3/2})$. On the other hand, by (d),
$\sum_{n\le x}r_4(n)=8\sum_{m\le x}\sum_{d\le x/m,\,4\nmid d}d$, and $\sum_{d\le y,\,4\nmid d}d=\frac{y^2}2-4\cdot\frac12\bigl(\frac y4\bigr)^2+O(y)=\frac38y^2+O(y)$. Hence
\[
\sum_{n\le x}r_4(n)=3x^2\sum_{m\le x}\frac1{m^2}+O\Bigl(x\sum_{m\le x}\frac1m\Bigr)=3\zeta(2)x^2+O(x\log x).
\]
Comparing the two leading terms, $3\zeta(2)=\frac{\pi^2}2$.

\tools{Jacobi's triple product and identity (Problem~\ref{pb:ramanujan5}); Lambert series; lattice points and volumes (Volume~1); $\sum_{m\le x}\frac1{m^2}=\zeta(2)+O(\frac1x)$.}

\leads Jacobi found the formula in the \emph{Fundamenta Nova} of 1829 with elliptic functions; the elementary route through $E^6$ is Hirschhorn's (1987). The two counts in (f)
give a proof of $\zeta(2)=\frac{\pi^2}6$ by counting points in a four-dimensional ball, with the volume $\frac{\pi^2}2R^4$ carrying the $\pi$. In modular language, $\Theta(\tau)^4$ is a modular form of weight $2$ for $\Gamma_0(4)$, a two-dimensional space spanned by Eisenstein
series, which is why a pure divisor sum comes out.
\end{solution}

\begin{solution}{pb:densities}
\idea Modulo a high power of $p$ a solution can be lifted coordinate by coordinate, as in Newton's method, provided one coordinate is a unit. So the number of solutions
modulo $p^k$ grows exactly like $p^{3k}$, and the constant of proportionality is read off modulo $p$ (or modulo $8$ when $p=2$).

\writeup Write $Q(x)=x_1^2+x_2^2+x_3^2+x_4^2$. (a) Let $x$ be a primitive solution modulo $p^k$, $k\ge1$, say $p\nmid x_1$, and $Q(x)=n+p^kc$. Its $p^4$ lifts
$x+p^kt$ modulo $p^{k+1}$ satisfy $Q(x+p^kt)\equiv n+p^k(c+2\sum x_it_i)\pmod{p^{k+1}}$, since $2k\ge k+1$. The condition $c+2\sum x_it_i\equiv0\pmod p$ is a linear equation in $t$
with the unit coefficient $2x_1$, so exactly $p^3$ lifts are solutions; every primitive solution modulo $p^{k+1}$ arises once in this way. A non-primitive solution is $x=py$,
with $y$ determined modulo $p^{k-1}$, and $Q(x)=p^2Q(y)$. If $p^2\nmid n$ there is none ($k\ge2$). If $p^2\mid n$, the condition is $Q(y)\equiv\frac n{p^2}\pmod{p^{k-2}}$, which
depends on $y$ modulo $p^{k-2}$; each such class has $p^4$ lifts modulo $p^{k-1}$, so there are $p^4N_{p^{k-2}}(\frac n{p^2})$ of them.

(b) By (a), $p^{-3k}N^*_{p^k}(n)=p^{-3}N^*_p(n)=:D^*$ for all $k\ge1$. By Problem~\ref{pb:sumsquares}, $N_p(n)=p^3-p$ if $p\nmid n$, and $N_p(n)=p^3+p^2-p$ if $p\mid n$; removing the
zero vector in the second case, $D^*=1-p^{-2}$ if $a=0$, and $D^*=p^{-3}(p^3+p^2-p-1)=(1-p^{-2})(1+p^{-1})$ if $a\ge1$. For $a\le1$ all solutions modulo $p^k$, $k\ge2$, are
primitive, so $\delta_p(n)=D^*$. For $a\ge2$, (a) gives $p^{-3k}N_{p^k}(n)=D^*+p^{-2}\,p^{-3(k-2)}N_{p^{k-2}}(\frac n{p^2})$, so by induction on $a$ the limit exists and
$\delta_p(n)=(1-p^{-2})(1+p^{-1})+p^{-2}(1-p^{-2})(1+\dots+p^{-(a-2)})=(1-p^{-2})(1+p^{-1}+\dots+p^{-a})$.

(c) Let $x$ be primitive modulo $2^k$, $k\ge3$, with $x_i$ odd. All $16$ lifts $x+2^kt$ have $Q\equiv Q(x)+2^{k+1}\sum x_it_i\equiv Q(x)\pmod{2^{k+1}}$, so either all or none are
solutions modulo $2^{k+1}$. Let $x'=x+2^{k-1}e_i$, where $i$ is the first odd coordinate: $Q(x')=Q(x)+2^kx_i+2^{2k-2}\equiv Q(x)+2^k\pmod{2^{k+1}}$ as $2k-2\ge k+1$. So
$x\mapsto x'$ is an involution without fixed points on the primitive solutions modulo $2^k$, and in each pair exactly one element has its $16$ lifts solutions. Hence
$N^*_{2^{k+1}}=\frac{16}2N^*_{2^k}$. Modulo $8$, odd squares are $1$ and even squares are $0$ (for $0,4$) or $4$ (for $2,6$). If $j$ coordinates are odd and $m$ of the even ones are
$\equiv2\pmod4$, then $Q\equiv j+4m\pmod8$. Counting, $N^*_8(n)=4\cdot4\cdot32=512$ for odd $n$ ($j\equiv n\pmod4$), $6\cdot16\cdot8=768$ for $n\equiv2\pmod4$ ($j=2$), $256$ for
$n\equiv4\pmod8$ ($j=4$), and $0$ for $8\mid n$. So $D^*=2^{-9}N^*_8(n)=1,\frac32,\frac12,0$ in these cases, and as in (b), $\delta_2(n)=D^*+\frac14\delta_2(\frac n4)$ when $4\mid n$. This
gives $\delta_2=1$ for $a=0$, $\frac32$ for $a=1$, $\frac12+\frac14=\frac34$ for $a=2$, and $\delta_2(n)=\frac14\delta_2(\frac n4)$ for $a\ge3$; by induction $\delta_2(n)=3\cdot2^{-a}$ for $a\ge1$.

(d) For odd $p\nmid n$, $\delta_p=1-p^{-2}$, and $\prod_{p>2}(1-p^{-2})=\frac{6/\pi^2}{3/4}=\frac8{\pi^2}$. Let $n=2^au$ with $u$ odd. Since $(1+p^{-1}+\dots+p^{-a_p})=\frac{\sigma(p^{a_p})}{p^{a_p}}$,
\[
\pi^2n\prod_p\delta_p(n)=\pi^2\cdot2^au\cdot\delta_2(n)\cdot\frac8{\pi^2}\cdot\frac{\sigma(u)}u=8\cdot2^a\delta_2(n)\,\sigma(u),
\]
which is $8\sigma(u)$ for $a=0$ and $24\sigma(u)$ for $a\ge1$: exactly $r_4(n)$ by Problem~\ref{pb:jacobi}(d). The product converges absolutely because all but finitely many factors are
$1-p^{-2}$. The ball $Q(x)\le t$ has volume $\frac{\pi^2}2t^2$, whose derivative at $t=n$ is $\pi^2n$: the real solutions of $Q(x)=n$ have ``density'' $\pi^2n$, just as
$p^{-3k}N_{p^k}(n)$ is the density of solutions modulo $p^k$.

\tools{Hensel's lifting (Volume~2); $N_p(n)$ from Problem~\ref{pb:sumsquares}; $\zeta(2)=\frac{\pi^2}6$; Problem~\ref{pb:jacobi}.}

\leads ``Number of solutions $=$ density at infinity $\times$ product of local densities'' is the Hardy--Littlewood circle method in its simplest form. For four squares the
formula is exact; for a positive definite form $Q$ in general, Siegel's mass formula (1935) says that it holds for a weighted average over the forms of the same genus, and
the circle method proves it asymptotically for five or more variables. Problem~\ref{pb:jacobi} computed the left side and this problem the right side, independently.
\end{solution}

\section*{Sources and further reading}

\begin{itemize}
\item T.~M.~Apostol, \emph{Introduction to Analytic Number Theory}, Springer, 1976. The construction of Dirichlet characters, the Euler product, and Dirichlet's theorem on primes in arithmetic progressions.
\item H.~Davenport, \emph{Multiplicative Number Theory}, 3rd ed., revised by H.~L.~Montgomery, Springer, 2000. The analytic theory built on the products of this chapter.
\item G.~Tenenbaum, \emph{Introduction to Analytic and Probabilistic Number Theory}, 3rd ed., American Mathematical Society, 2015. Sieves beginning from Legendre's formula, and the reasons an exact formula is not yet an asymptotic one.
\item Volume~6 of this series proves Dirichlet's theorem, the prime number theorem and the first modern sieves. Volume~13 identifies the characters that describe prime splitting. Volume~14 returns to multiplicative functions and to the limitations of sieve weights.
\end{itemize}
